\section{Conclusion}
\label{sec:conclusion}

This work demonstrates that CatBoost trained on approximately one million Korean NHIS health checkup records with 40 routinely collected features achieves an ROC-AUC of 0.817 for T2D screening, matching neural network architectures while providing native SHAP interpretability. Three-tier feature ablation reveals that strong screening performance persists without lipid panel data (Core+Labs AUC 0.808), with the full lipid panel contributing only a marginal gain of +0.009, an important finding given that lipid tests are optional in many screening programs. At the chosen clinical threshold the model attains 85.2\% sensitivity and 98.0\% negative predictive value, supporting its suitability for population-level triage. SHAP analysis confirms that the most influential predictors---age, liver enzymes, waist circumference, and lipid ratios---align with established clinical risk factors, reinforcing model trustworthiness. The deployed web-based screening tool translates calibrated probabilities into four interpretable risk tiers, requiring no imputation and handling missing data natively. Future work should prioritize multi-national validation to assess generalizability beyond the Korean population and to evaluate performance across diverse healthcare systems.
